% =====================================================================
% Module 02 — Fiber Bundle Structure
% Academic-Galactic v1.2
% =====================================================================

\section{Fiber Bundle Structure}
\label{sec:FiberBundleStructure}

\subsection{Overview}
The fiber bundle structure of $\mathcal{M}^{80}$ provides the geometric
connective tissue between filtration layers and operator-induced transitions.
Each filtration layer $\mathcal{F}_k$ carries a smooth fiber bundle
\[
\pi_k : E_k \to \mathcal{F}_k,
\]
where $E_k$ encodes curvature envelopes, holonomy transport, and local geometric
invariants.

Curvature and holonomy foundations follow classical differential geometry
\citeClassA{Atiyah1956} and modern holonomy-curvature research
\citeClassA{BlackholeHolonomy2026,GutierrezUle2026,LiPengDong2026}.

\subsection{Curvature Envelopes}
\begin{definition}[Curvature Envelope]
A curvature envelope on $\mathcal{F}_k$ is a smooth assignment
\[
\Phi_7(p) \in \Omega^2(T_p\mathcal{F}_k)
\]
representing the local curvature induced by the operator altitude.
\end{definition}

Curvature envelopes satisfy spectral-nonnegativity constraints
\citeClassA{SpectralCurvature2026} and interact with collapse-surface
stratification (Module 04).

\subsection{Holonomy Transport}
Holonomy operators $\Phi_8$ define parallel transport along curves
$\gamma : [0,1] \to \mathcal{F}_k$:
\[
\Phi_8(\gamma) : E_{\gamma(0)} \to E_{\gamma(1)}.
\]

Holonomy behavior is grounded in:
\begin{itemize}
    \item classical holonomy theory \citeClassA{Atiyah1956},
    \item curvature-holonomy interaction \citeClassA{BlackholeHolonomy2026},
    \item informational holonomy curvature \citeClassA{GutierrezUle2026},
    \item discrete-to-continuous holonomy convergence \citeClassA{LiPengDong2026}.
\end{itemize}

\subsection{Holonomy–Curvature Compatibility}
Holonomy transport satisfies the curvature compatibility condition:
\[
\Phi_8(\gamma) = \exp\left( \int_\gamma \Phi_7 \right),
\]
ensuring that holonomy is generated by curvature envelopes.

\subsection{Fiber Bundle Local Triviality}
Each bundle $E_k$ is locally trivial:
\[
\pi_k^{-1}(U) \cong U \times F_k,
\]
where $F_k$ is the fiber type at altitude $k$.

Local triviality ensures compatibility with:
\begin{itemize}
    \item spectral bundle structure (Module 03),
    \item collapse-surface stratification (Module 04),
    \item operator algebra domain constraints (Module 05).
\end{itemize}

\subsection{Holonomy-Induced Drift}
Holonomy operators induce drift along collapse surfaces:
\[
\Phi_8^t : \mathcal{F}_k \to \mathcal{F}_k,
\]
where $\Phi_8^t$ is the topology-modulation operator.

This drift is responsible for:
\begin{itemize}
    \item collapse-surface deformation,
    \item topology modulation,
    \item stratified surface evolution.
\end{itemize}

\subsection{Compatibility with Filtration}
Fiber bundle structure is filtration-preserving:
\[
i_k^\ast(E_{k+1}) = E_k.
\]
This ensures that curvature envelopes and holonomy transport remain consistent
across altitude transitions.

\subsection{Summary}
The fiber bundle structure defines the geometric backbone of $\mathcal{M}^{80}$,
including curvature envelopes, holonomy transport, topology modulation, and
collapse-surface drift. These structures interact directly with the operator
altitude and provide the geometric substrate for spectral continuity and
operator algebra behavior.

% =====================================================================
% End Module 02
% =====================================================================
